\clearpage

\section{Integrals}

What is it other than the area under curve?

Interesting examples where the Riemann's integration fails is that
$f(x) = 1$ if $x$ is rational and $0$ otherwise over the interval $[0,1]$.
Why this function's integral is problematic in Riemann's definition.

Lebesgue used measure theory to define integral.

\begin{quote}
	Many times, when Riemann fails, Lebesgue works.
\end{quote}

Set being countable and uncountable rely on if we can arrange them in
a `readable' sequence, could be infinite like rational numbers.

Proving that $\Rf$ is uncountable is not so trivial: decimal expansion.

No matter how hard you try to write the numbers in $[0,1)$,
there is always some values left out.

Continuum hypothesis: How do you check which $\infty$ is bigger?
Or which set has the higher cardinality?

Definition 1.2.2 ()

\begin{example}[Set of Lebesgue measure zero, 1.2.2]
	A set $S$ is zero in Lebesgue measure
	when it can be covered with a sequence of open intervals ($I_1, I_2, \dots$).
	And the sum $\sum_{0}^{\infty} m(I_n)$ can be made arbitrarily small
\end{example}

In other words, for any $\epsilon > 0$, we can always find a sequence of $I_n$
that covers $S$.

\begin{thm}
	Any countable infinite set of $S$ has Lebesgue measure zero.
\end{thm}

\begin{remark}
	The set of measure zero would not change a damn thing on $f$
\end{remark}

\subsection{Filtrations and sigma-fields}

\begin{define}[Filtration]
	Let $(\Omega, \Ac, P)$ be a probability space. An increasing sequence $\Fc_n$ of sub-$\sigma$ algebras of $\Ac$
	(i.e. $\Fc_0 \subseteq \Fc_1 \subseteq \cdots \Fc_n \subseteq \cdots \subseteq \Ac$) is called a filtration.
\end{define}

The sub-$\sigma$-algebra is just the sigma algebra of the $X_0, X_1, X_n$.

What is the difference between the stopping time and random times that can take possibly infinity.

Filtration sigma-field at time $t$, filtration is a n increasing sequence of sigma-fields.

What is the lim sup here? $A_t$ is an adapted sequences of events in some filtered probability space.
$$
	A_\infty := \limsup_{t\to \infty} A_t := \bigcap_{t\in\Nf}\bigcup_{s\geq t}A_s
$$

\begin{define}[Absolute continuous]
	Let $p, q$ be two probability distributions. $p$ is called \emph{absolutely continuous} with respect to $q$ if the

	More explicitly, this reads ($p \ll q$) that
	\begin{align*}
		q(A) = 0 \quad \Rightarrow \quad p(A) = 0,  \quad \text{for any } A \in \Fc_t \text{ and } t \in \Nf.
	\end{align*}
\end{define}

What about \emph{CORASENING}?
